\begin{center}
\setstretch{1.3}
\sanhao\bf Study on the internal structure of exotic hadronic states from the perspective of hadron correlation functions and invariant mass distributions
\end{center}
\vspace{0.7cm}
Graduate student: Hai-Peng Li\\
Supervisor: Prof. Wei-Hong Liang\\
%Graduate student: $\square\square\square$\\
%Supervisor: Prof. $\square\square\square$\\
Major: Physics\\
Research direction: Particle Physics\\
Grade: 2022\\
\vspace{0.5cm}
\centerline{\bf\xiaosanhao Abstract}
\vspace{-1cm}
\par

The structure and properties of exotic hadronic states have been a longstanding topic of interest in hadron physics.
With the increasing number of exotic hadron candidates being discovered in experiments, it has become particularly important to develop effective methods to investigate their nature.
Recently, the ALICE and STAR collaborations have reported new experimental data about femtoscopic correlation functions,
from which scattering lengths and effective ranges have been extracted.
These developments offer a novel avenue for studying exotic hadronic states.

This thesis explores the potential of using scattering lengths and effective ranges to probe the nature of exotic hadronic states.
First, within the framework of Chiral Unitary Approach, we study the resonance $N^{*}(1535)$, and calculate its pole position, scattering lengths, effective ranges and correlation functions.
These observables are treated as pseudo-data from experiments and analyzed under the genuine picture.
The determined observables differ significantly from those predicted by the molecular picture, demonstrating that such observables can effectively distinguish between different structural interpretations.

Building on this, we further investigate whether these observables can be accurately determined from correlation functions, and study exotic hadronic states in both the light and heavy sectors.
We focus on the exotic candidate $\Sigma^{*}(1430)$ in light sector, located near the $\bar{K}N$ threshold.
Within the Chiral Unitary Approach framework, we calculate correlation functions for relevant channels.
By employing a method as minimal model dependent and resampling, we successfully determine the scattering lengths and effective ranges with uncertainties below 20\%, demonstrating the feasibility of high precision determination from correlation functions.
Based on this success, we extend the method to the heavy sector, specifically the $BD$ interaction.
Although the bound state mass lies below the threshold and correlation functions provide only information from its tail, we still manage to determine its properties: the binding energy with approximately 50\% uncertainty, and the molecular probability with 6\% precision.
The scattering lengths are determined with good accuracy, while the effective ranges show larger uncertainty but remains consistent with theoretical predictions.
These results confirm the applicability of correlation functions to heavy flavor exotic hadronic state.
Hece, the determination of observables from correlation functions is shown to be a generally valid method for probing the properties of exotic hadronic states.

We also examine the possibility of determining observables from invariant mass distributions, including the cusp from virtual state and near threshold distributions of bound state.
For the long-debated $K^{+}\bar{K}^{0}$ scattering length, we analyzed the invariant mass distributions in the decay $\chi_{c1}\to \eta\pi^{+}\pi^{-}$, and successfully determined the $K^{+}\bar{K}^{0}$ scattering length and effective range, with uncertainties below 20\%.
Moreover, by analyzing the invariant mass distributions with 50 MeV above $\bar{D}\bar{K}$ threshold in the decay $\Lambda_{b}\to \Lambda_{c}^{+}\bar{D}^{0}K^{-}(D^{-}\bar{K}^{0})$, we successfully determine the mass of the $D^{*}_{s0}(2317)$ with a 4 MeV uncertainty and determine its molecular probability with an 11\% uncertainty.

In summary, this thesis demonstrates the crucial role of scattering observables in exploring the nature of exotic hadrons.
These quantities can be determined not only from correlation functions but also, in certain cases, from invariant mass distributions with high precision.
This work presented here provides a theoretical foundation for future experimental studies, and with ongoing experimental advancements, we can expect to gain further insights into the nature of exotic hadrons.

%%%%%%%%%%%

\vspace{0.30cm}
\noindent{\bf Keywords: } {Strong interaction; Hadronic structure; Exotic hadrons; Correlation functions; invariant mass distributions}
%强相互作用；强子结构；奇特强子态；关联函数；不变质量分布